\section{问题重述、数据分层与符号}
\subsection{任务与递进关系}
题目与随附数据说明\cite{problem}要求依次回答四项任务：从多指标评价预训练文本质量并分析冲突、建立领域配比与训练损失的关系；结合模型参数量 $N$、训练数据量 $D$、质量 $Q$ 与配比 $p$ 讨论标度律和边际效用；在给定算力预算下配置资源；最后分析开放模型能力变化和算力增速放缓的情景。本文以题目附件为主证据。四问按“质量和配比候选$\rightarrow$标度响应$\rightarrow$预算配置$\rightarrow$能力边界”的顺序连接，但只对有共同观测协议的量作数值传递。

\subsection{数据角色及不可混用的标尺}
A1--A3 是质量指标记录，A4--A15 提供配比及不同规模的 Loss，A16 为领域映射参考；B1 为训练检查点，B3 是插值轨迹，B6--B8 涉及半合成质量实验；C1--C8 包含开放模型排行榜、权重与许可、算力及任务信息。B3 插值记录不用于独立验证；A12--A15 的大尺度 Loss 估计不当成真实训练观测。$Q_A$ 表示问题一的相对质量评分，$Q_B$ 表示 B6 的原生质量水平。两者没有相同样本的成对校准数据，不能直接互换。不同来源的 Loss 也不强行拼接为同一训练集。

本文令 $N_B=N/10^9$、$D_B=D/10^9$，即参数量和 token 数均以十亿计；$C$ 以 FLOPs 计。$p=(p_1,\ldots,p_{17})$ 为非负且和为 1 的领域配比。$L$ 表示各附件自身协议下的 Loss；排行榜六任务均分以分数点计，与 $L$ 不同量纲。

\begin{table}[htbp]\centering\caption{四问主要证据及可回答范围}
\small\begin{tabular}{p{1.3cm}p{3.1cm}p{7.3cm}}\toprule
问题 & 主证据 & 可回答范围\\\midrule
一 & A1--A7、A16 & 数据内相对质量、1M 配比预测与候选配方\\
二 & B1、B6；B4/B5/B7 验证 & 标度关系、半合成质量效应和条件边际量\\
三 & 第一问候选配比、第二问响应、题设成本 & 指定支持域与成本函数下的条件配置\\
四 & C1/C2/C4/C6/C8 & 开放样本的描述性趋势、算力情景与预测可识别性\\\bottomrule
\end{tabular}\end{table}

\subsection{建模原则}
先核对文件编号、行数与字段，再划分训练、验证、插值和估算数据；同一模型家族或训练轨迹整体留出，避免相邻记录泄漏。所有模型结论附上支持域、误差口径和证据类型。题目允许对配比采用前问结果作为固定方案，本文在第三问采用这一选择，并明确其尚未经跨问联合实验检验。
